% authority: science_draft
\documentclass[11pt]{article}

\usepackage[margin=1in]{geometry}
\usepackage{amsmath,amssymb,amsthm}
\usepackage{booktabs}
\usepackage{longtable}
\usepackage{enumitem}

\newcommand{\EqMScert}{\mathrm{EqMS}^{\mathrm{cert}}_{\mathrm{src}}}
\newcommand{\NarrowMSCertEq}{\mathrm{NarrowMSCertEq}_{v1}}
\newcommand{\CertSrc}{\mathsf{Cert}_{\mathrm{src}}}
\newcommand{\CTr}{C_{\mathrm{transport}}^{\mathrm{src}}}
\newcommand{\CInv}{C_{\mathrm{invariance}}^{\mathrm{src}}}
\newcommand{\CFact}{C_{\mathrm{factorization}}^{\mathrm{src}}}
\newcommand{\CNTI}{C_{\mathrm{no\text{-}target}}^{\mathrm{src}}}
\newcommand{\DeclEq}{\equiv_{\mathrm{decl}}^{\mathrm{src}}}
\newcommand{\Bottom}{\bot}
\newcommand{\RRELaw}{\mathrm{RR}_{E}\mathrm{TransportCompletenessOrInvarianceLaw}_{v1}}
\newcommand{\PositiveMSProfile}{\mathrm{PositiveMSProfile}_{v1}}
\newcommand{\SMSAR}{\mathrm{SourceMatterSemanticsAdoptionReadinessLaw}_{v1}}
\newcommand{\MetricData}{\mathrm{MetricData}}
\newcommand{\geff}{g_{\mathrm{eff}}}

\newtheorem{finding}{Finding}
\newtheorem{obligation}{Obligation}
\newtheorem{boundary}{Boundary}
\newtheorem{nonconclusion}{Non-conclusion}

\title{EqMS Definition/Theorem-Content Separation Audit v16}
\author{Ontology Formalizer Packet RT-20260705-004}
\date{2026-07-05}

\begin{document}
\maketitle

\section{Control Status}

This artifact implements v16 P5-T01. It audits \(\NarrowMSCertEq\) and the
source certificate algebra to separate definition unfolding from theorem
content. The audit is draft/control. It clarifies scope; it does not demote
the existing scoped evidence-status and does not promote any downstream
physics claim.

This packet does not edit canonical ontology TeX. It does not adopt a source
law, \(\RRELaw\), \(\PositiveMSProfile\), \(\SMSAR\), matter semantics,
detector semantics, a coupling law, matter coupling, \(\MetricData(E)\),
scoped \(\geff\), stress-energy semantics, a stress-energy tensor, matter
action, Einstein equations, benchmark status, or completed derivation.

\section{Source Definitions Imported From P2/P3}

\begin{longtable}{p{0.24\linewidth}p{0.18\linewidth}p{0.42\linewidth}}
\toprule
Imported source item & Classification & Audit result \\
\midrule
\(\EqMScert\) predicate & \texttt{validity\_predicate} &
The predicate is the conjunction of at least one explicit source certificate,
\(\CNTI\), declared-object consistency, and inactive fail-closed branches. \\
\(\CTr,\CInv,\CFact\) certificate classes & \texttt{constructor\_definition} &
They define allowed source-side witness kinds. Existence of such witnesses is
not automatic. \\
\(\CNTI\) no-target guard & \texttt{validity\_predicate} &
It is a negative hygiene condition. It is not a positive transport,
invariance, or factorization certificate. \\
Missing, malformed, target-import, detector-import, and process-authority
branches & \texttt{countermodel\_or\_obstruction} &
They explain when attempted equivalence returns \(\Bottom\) or records a
certificate-gap obstruction. \\
Identity, composition, and restriction primitives & \texttt{constructor\_definition} &
The primitives define expressions whose laws require separate conditional
proofs. \\
\bottomrule
\end{longtable}

\section{Theorem Statements Imported From P2/P3}

\begin{longtable}{p{0.25\linewidth}p{0.20\linewidth}p{0.39\linewidth}}
\toprule
Statement & Classification & Audit result \\
\midrule
\(\NarrowMSCertEq\): explicit certificates imply \(A_{\mathrm{src}}
\EqMScert B_{\mathrm{src}}\) & \texttt{definition\_unfolding\_theorem} &
The positive proof mainly checks the clauses in the definition of
\(\EqMScert\). Its value is boundary discipline, not an independent
certificate-existence theorem. \\
Certificate-gap fail-closed proposition & \texttt{countermodel\_or\_obstruction} &
If no explicit certificate exists, the theorem cannot be used; scoped evidence
or process status cannot fill the slot. \\
Malformed/imported certificate fail-closed proposition &
\texttt{countermodel\_or\_obstruction} &
Invalid witnesses cannot establish \(\EqMScert\). \\
Identity declared-scope preservation & \texttt{nontrivial\_conditional\_theorem} &
The law proves preservation inside declared scope under primitive
well-formedness assumptions. \\
Source-side closure under compatible composition &
\texttt{nontrivial\_conditional\_theorem} &
The law adds conditional closure for compatible valid certificate records. \\
Source restriction closure & \texttt{nontrivial\_conditional\_theorem} &
The law proves closure only for declared source subdomains and fails closed for
non-source subdomains. \\
Malformed, missing, and target-import fail-closed laws &
\texttt{countermodel\_or\_obstruction} &
The laws locate precise failure branches and block adoption-style overread. \\
\bottomrule
\end{longtable}

\section{Proof Steps That Are Definitional Unfoldings}

\begin{finding}[Definition-unfolding part of \(\NarrowMSCertEq\)]
The proof of the positive \(\NarrowMSCertEq\) branch unfolds the definition of
\(\EqMScert\): one of \(\CTr,\CInv,\CFact\) is assumed; \(\CNTI\) is assumed
to pass; declared-object indices are assumed consistent; fail-closed branches
are assumed inactive. Under those premises, the conclusion is exactly the
predicate defined by those clauses.
\end{finding}

The transport, invariance, and factorization cases in the P2 theorem therefore
classify as \texttt{definition\_unfolding\_theorem} unless a packet supplies
new independent proof that the relevant certificate exists or that a family of
certificates is closed under an operation.

\section{Proof Steps That Are Nontrivial}

\begin{longtable}{p{0.28\linewidth}p{0.18\linewidth}p{0.38\linewidth}}
\toprule
Source-side content & Classification & Why nontrivial \\
\midrule
P3 identity declared-scope preservation &
\texttt{nontrivial\_conditional\_theorem} &
It proves an operation law for valid identity certificates and compatible
composition inside scope intersections. \\
P3 compatible composition closure &
\texttt{nontrivial\_conditional\_theorem} &
It checks domain-codomain compatibility, source-only inheritance, and endpoint
preservation for composed certificates. \\
P3 restriction closure & \texttt{nontrivial\_conditional\_theorem} &
It separates source subdomains from target, detector, benchmark, and process
substitutes. \\
P4 transport instance \texttt{SCI-TRANSPORT-001} &
\texttt{finite\_instance\_theorem} &
A finite table verifies row-wise source transport of labels, guards, and
response tokens. \\
P4 invariance instance \texttt{SCI-INVARIANCE-001} &
\texttt{finite\_instance\_theorem} &
The non-identity cycle \(\sigma=(0\,1\,2)\) preserves the declared successor
relation by direct enumeration. \\
P4 factorization instance \texttt{SCI-FACTORIZATION-001} &
\texttt{finite\_instance\_theorem} &
The equality \(h=q\circ p\) is checked on every finite source-domain row. \\
P4 negative instances & \texttt{countermodel\_or\_obstruction} &
They show explicit ways the predicate fails closed when a witness is absent,
malformed, target-importing, detector-importing, stress-energy based, or
process-status based. \\
\bottomrule
\end{longtable}

\section{Missing Theorem Content}

The exact theorem-content gaps are:

\begin{enumerate}[label=\textbf{G\arabic*.},leftmargin=*]
  \item No general certificate-existence theorem: nothing proves that a broad
  class of source matter-semantics objects admits valid \(\CTr,\CInv\), or
  \(\CFact\) witnesses.
  \item No inverse or symmetry theorem for arbitrary valid source
  certificates. Identity and compatible composition are not enough to prove an
  equivalence relation in the usual algebraic sense.
  \item No theorem that scoped evidence/preconditions such as
  \(\PositiveMSProfile\), \(\SMSAR\), or \(\RRELaw\) instantiate concrete
  certificate bundles.
  \item No theorem that finite/local P4 certificate instances generalize to
  unrestricted source families.
  \item No theorem connecting \(\EqMScert\) to detector semantics,
  physical matter semantics, stress-energy semantics, matter action, coupling
  law, matter coupling, \(\MetricData(E)\), \(\geff\), or Einstein equations.
  \item No theorem that no-target hygiene alone supplies positive semantics.
\end{enumerate}

\section{Future Theorem Obligations}

\begin{longtable}{p{0.26\linewidth}p{0.20\linewidth}p{0.38\linewidth}}
\toprule
Future target & Classification & Required burden \\
\midrule
Certificate existence for a named finite/local source-object family &
\texttt{future\_theorem\_target} &
State a domain of source objects and construct at least one explicit
certificate family or prove a precise obstruction. \\
Inverse/symmetry condition for declared certificates &
\texttt{future\_theorem\_target} &
Define when a valid certificate has a source-side inverse and when it does not. \\
Associativity or category-like closure for compatible composition &
\texttt{future\_theorem\_target} &
Prove the law with explicit source scopes, or record the obstruction. \\
Finite-to-family lift for P4 instances &
\texttt{future\_theorem\_target} &
Specify the finite/local generalization parameter and preservation conditions. \\
Adoption-readiness separation theorem &
\texttt{future\_theorem\_target} &
Prove that scoped evidence can remain evidence/precondition without becoming
adopted semantics, source law, or coupling law. \\
\bottomrule
\end{longtable}

\section{Finite/Local Examples From P4}

P4 demonstrates why theorem content is not exhausted by naming the
\(\EqMScert\) predicate.

\begin{itemize}[leftmargin=*]
  \item \texttt{SCI-TRANSPORT-001} supplies an explicit transport table on
  \(\{0,1,2\}\). The nontrivial work is the row-wise source check.
  \item \texttt{SCI-INVARIANCE-001} supplies a non-identity cyclic relabeling.
  The nontrivial work is the direct proof that the successor relation is
  preserved.
  \item \texttt{SCI-FACTORIZATION-001} supplies a middle object and maps
  \(p,q,h\). The nontrivial work is the direct proof of \(h=q\circ p\).
  \item The negative P4 packet supplies eight countermodel or obstruction
  fixtures. These fixtures prove that absent payloads, malformed maps, changed
  factor objects, target metrics, detector semantics, stress-energy shortcuts,
  validator passes, and scoped evidence cannot substitute for explicit source
  certificates.
\end{itemize}

\section{Forbidden Conclusion Boundary}

\begin{boundary}[Audit boundary]
This audit preserves \(\NarrowMSCertEq\) as scoped source-extension
evidence-status for its declared conditional theorem role. It clarifies that
the positive proof is largely definition unfolding and that future theorem
work must target certificate existence, operation closure, inverse/symmetry,
finite-family lifting, or precise obstruction. It does not revoke the scoped
status and does not promote that status into adoption.
\end{boundary}

\begin{nonconclusion}
No source-law adoption, \(\RRELaw\) adoption, unrestricted \(RR_E\) theorem,
\(\PositiveMSProfile\) adoption, \(\SMSAR\) law adoption, matter-semantics
adoption, detector-semantics adoption, coupling-law adoption, matter-coupling
derivation or adoption, stress-energy semantics, stress-energy tensor, matter
action, Einstein equations, benchmark promotion, or completed derivation is
authorized by this audit.
\end{nonconclusion}

\section{Machine-Readable Summary}

\begin{verbatim}
audit_id: "eqms_definition_theorem_content_separation_audit_v16"
implemented_plan_task: "P5-T01"
primary_result:
  NarrowMSCertEq_v1_positive_branch:
    classification: "definition_unfolding_theorem"
    reason: "positive conclusion is the EqMS_src^cert predicate under assumed explicit certificates"
  certificate_operation_laws:
    classification: "nontrivial_conditional_theorem"
  p4_positive_instances:
    classification: "finite_instance_theorem"
  p4_negative_instances:
    classification: "countermodel_or_obstruction"
theorem_content_gap_ids:
  - "G1 general_certificate_existence"
  - "G2 inverse_or_symmetry"
  - "G3 scoped_evidence_to_certificate_bundle"
  - "G4 finite_to_family_lift"
  - "G5 physical_semantics_or_coupling_link"
  - "G6 no_target_guard_as_positive_semantics"
next_route: "P5-T02 refactored source-equivalence target specification"
physics_promotion_authorized: false
\end{verbatim}

\section{APA 7 Source Materials}

The AEther-Flow Research Project. (2026, July 2). \textit{Source-side
matter-semantics object and certificate manifest v1} [Research-control TeX
artifact].

The AEther-Flow Research Project. (2026, July 2). \textit{Narrow source-side
matter-semantics equivalence theorem v1} [Research-control TeX artifact].

The AEther-Flow Research Project. (2026, July 2). \textit{Matter-semantics
equivalence theorem Refuter stress v1} [Research-control TeX artifact].

The AEther-Flow Research Project. (2026, July 2). \textit{NarrowMSCertEq v1
scoped evidence-status Gate Chair review} [Research-control TeX artifact].

The AEther-Flow Research Project. (2026, July 2). \textit{Source certificate
algebra primitives v1} [Research-control TeX artifact].

The AEther-Flow Research Project. (2026, July 2). \textit{Source certificate
operation laws and fail-closed lemma v1} [Research-control TeX artifact].

The AEther-Flow Research Project. (2026, July 4). \textit{Finite local valid
transport certificate instance v1} [Research-control TeX artifact].

The AEther-Flow Research Project. (2026, July 4). \textit{Finite local valid
invariance certificate instance v1} [Research-control TeX artifact].

The AEther-Flow Research Project. (2026, July 5). \textit{Finite local valid
factorization certificate instance v1} [Research-control TeX artifact].

The AEther-Flow Research Project. (2026, July 5). \textit{Negative certificate
instance packet v1} [Research-control TeX artifact].

\end{document}
